\chapter{What follows mathematically, and what does not}\label{ch:limits}
\question{Which conclusions follow from definitions, and which need more evidence?}

A proof is only as broad as its assumptions. This chapter collects the exact conclusions and their limits. These are mathematical statements about a declared model; its physical adequacy and institutional use need separate justification.

\section{A small theorem ledger}
\textit{Affected-support cancellation.} Under one fixed additive potential, unchanged terms cancel. The local account must include every affected factor, including coupled factors (Chapter~\ref{ch:potential}).

\textit{Exact quadratic expansion.} Fixed references and positive scales give equation~\eqref{eq:finite}. Its positive correction for a nonzero increment makes the initial tangent overstate finite field reduction.

\textit{Endpoint/path equality.} For a fixed continuously differentiable potential on a declared differentiable path, integrating its rate of change gives its endpoint difference. Equation~\eqref{eq:path} assumes fixed additive increments; it does not establish a physical history.

\textit{Common-path closure.} Under those assumptions, the $R_a$ lines sum to the group field value. Closure does not select owners or establish fairness.

The scoped historical foundation and Gaussian reconciliation support these results.\cite{foundation,gaussian} We now derive a sequence identity for complete joint transitions.

\section{How the middle states cancel}
Suppose a first authoritative transition takes $x^0$ to $x^1$, then a second takes $x^1$ to $x^2$. Superscripts here label positions in a sequence, not powers. Use the same potential $V$ throughout and no unrecorded external changes between these transitions. With nonnegative process burdens $C_0$ and $C_1$, the two values are
\begin{align*}
 E_0&=V(x^0)-V(x^1)-C_0,\\
 E_1&=V(x^1)-V(x^2)-C_1.
\end{align*}
Adding cancels the two occurrences of $V(x^1)$, one negative and one positive. For $N$ consecutive transitions, the same cancellation gives
\begin{equation}\label{eq:telescoping}
 \sum_{n=0}^{N-1}E_n=V(x^0)-V(x^N)-\sum_{n=0}^{N-1}C_n.
\end{equation}
Every successor must be the next predecessor. Every transition is accounted once, under the same ruler and complete relevant state. A transition here may be the coherent joint action of a simultaneous group. This is a directly stated endpoint identity; it does not silently rewrite the one-action restrictions of the historical foundation's theorem.\cite{bridge,foundation}

In the two-tank illustration, a first three-litre delivery changes potential from four to $0.25$, giving $3.75$. A later one-litre delivery from that endpoint reaches $(10,10)$ and gives $0.25$. The sum is four. Quoting the later litre as if the first delivery had never happened would reuse the wrong baseline.

\section{Closed cycles and splitting}
If the complete endpoint returns to the complete start, $x^N=x^0$, the potential difference vanishes. Equation~\eqref{eq:telescoping} then gives
\begin{equation}\label{eq:cycle}
 \sum_n E_n=-\sum_n C_n\leq0.
\end{equation}
A genuine cost-free cycle has zero total field value. Positive included process burdens make the net negative. This is a conditional accounting conclusion about a closed undriven cycle with fixed $V$, complete state, faithful quantities and no duplication.

For a hand calculation, move three litres from $(14,6)$ to $(11,9)$, then move the same three litres back. The field values are $+3.75$ and $-3.75$. Their sum is zero. The calculation does not assert that physically moving real water back and forth has no resource cost; it is within the declared ideal lossless world.

Splitting a fixed endpoint change into consecutive pieces leaves the field total unchanged under these assumptions. Process burdens may differ if the actual split requires extra activations. A sequential field identity is not permission to ignore such burdens or to sum duplicated same-baseline singleton quotes.

Returning one visible stock is not necessarily a complete cycle. Fuel consumption, damage or an altered equipment state could remain elsewhere in a wider account. A cycle claim must restore the complete represented state relevant to it, not a convenient projection.

\section{External changes and a changed ruler}
Let external drive take $x^n$ to $z^n$ before a joint action takes $z^n$ to $x^{n+1}$. Define its signed potential change as
\begin{equation*}
 D_{{\rm ext},n}=V(z^n)-V(x^n).
\end{equation*}
Adding and subtracting $V(x^n)$ in each action value gives
\begin{equation}\label{eq:driven}
 \sum_n E_n=V(x^0)-V(x^N)
 +\sum_n D_{{\rm ext},n}-\sum_n C_n.
\end{equation}
The drive term has the same account unit as $V$. It can be positive or negative. A disturbance that brings the state closer to reference has a negative potential injection. The signed external record is not automatically an actor reward.\cite{foundation}

A different issue arises if the reference or scales change. Then adjacent terms may be evaluated under different versions of $V$ and no longer cancel. The model-change effect needs its own record. Rewriting earlier receipts under a new reference would hide that difference rather than resolve it.

These qualifications explain why the cycle identity is not a universal anti-gaming theorem. It does not settle fraud, omitted effects, coalition strategy, changed baselines or all driven histories. Historical adverse findings remain evidence about their original models. They are not erased by a cleaner exposition.

\section{A short prospective capacity bridge}
The Gaussian programme proposes persistent nonnegative experimental balances $B_i$ assigned to declared owners. These balances would constrain which actions a separately specified actor can afford. They are not physical water, installed pump capacity, money or a person's entitlement to care.

For the ideal $C=0$ proposal, suppose signed balance changes sum exactly to the joint action field reduction:
\begin{equation*}
 \sum_i\Delta B_{i,n}=V(z^n)-V(x^{n+1}).
\end{equation*}
Let an external audit $J$ accumulate the drive term, so $J_{n+1}=J_n+D_{{\rm ext},n}$. Add these declared updates to obtain
\begin{equation}\label{eq:capacity}
 V(x^{n+1})+\sum_i B_{i,n+1}-J_{n+1}
 =V(x^n)+\sum_i B_{i,n}-J_n.
\end{equation}
All terms use compatible dimensionless account units. The constant across events follows by cancellation. It depends on the stated update rules; it does not prove those rules desirable. Ownership, same-event netting and affordability are policy definitions beyond the identity. The current proposal introduces no borrowing, refill or freely transferable global pool.\cite{gaussian}

External drive changes the physical state and its audit; it does not itself earn actor capacity. A fixed audit after one disturbance still permits cycling, trapping or refusal. Under continuing drive, the physical state might have one kind of statistical behaviour while the balance and audit variables drift. A claim about the whole process would have to address them all.

\section{Why boundedness proves so little here}
In a closed nonnegative world with finite mass $M$, every stock satisfies $0\leq x_i\leq M$. Fixed finite references and positive scales therefore imply
\begin{equation}\label{eq:bound}
 V(x)\leq\frac12\sum_i
 \frac{\max\{(x_i^\star)^2,(M-x_i^\star)^2\}}{\sigma_i^2}<\infty.
\end{equation}
For each cell, the largest squared displacement on the interval occurs at one of its endpoints. Adding those bounds gives the displayed result, which may be loose because not all endpoint choices are jointly possible. It is still enough to establish boundedness.\cite{gaussian}

That bound holds before a beneficial EBU effect has been shown. A comparison policy in the same physical domain inherits it too. The interesting research questions concern recovery, time away from reference, repeated refusal, distribution and exposure, not unbounded growth of this fixed-domain potential.

\practice{A future study reports that both its EBU and comparison arms remained bounded. Has affordability been shown to cause stability?}{No. The finite-mass domain already supplies a bound. A stability, attraction or recovery claim requires the relevant definition, comparator, assumptions and evidence. An accounting invariant also does not supply the missing dynamics.}

\carry{The identities are useful because their meaning is precise. They bring us to the beginning of a research programme, not the end of one.}
